\appendix

\section{Establishing Lower Bounds for Inconsistency Rates}
\label{sec:inconsistency-calculation}

To determine an appropriate sample size for estimating the proportion of inconsistencies in Wikipedia, we compute the minimum number of claims required to achieve a 99\% confidence level with a 5\% margin of error using the Cochran formula~\citep{cochran1953sampling}:
\begin{align*}
n &= \frac{z^2 \times p(1-p)}{E^2}  \\ &= \frac{2.576^2 \times 0.5 \times 0.5}{0.05^2} = 664
\end{align*}
where $z = 2.576$ corresponds to a 99\% confidence level, $p = 0.5$ assumes maximum variance (yielding the largest required sample size), and $E = 0.05$ is the desired margin of error. This calculation indicates that we must examine at least 664 claims. We apply the same statistical approach to determine sample sizes for our analyses of the AmbigQA and FEVEROUS datasets.



\subsection{Inconsistencies Per Wikipedia Article Category}
\label{section:inconsistency-in-categories}
\begin{figure}[ht!]
\centering
\includegraphics[width=\linewidth]{assets/category_distribution_incocsistency.pdf}
\caption{Distribution of inconsistencies in Wikipedia across topics.}
\label{fig:category_inconsistency}
\end{figure}

The distribution of potentially inconsistent claims across Wikipedia categories reveals notable patterns (Figure~\ref{fig:category_inconsistency}). History exhibits the highest rate of inconsistencies (17.7\%), followed by Everyday Life (16.9\%) and Society \& Social Sciences (14.3\%). These trends are reflected in concrete cases from our analysis. For example, the claim that ``The Ottoman Empire first developed the technique of using explosive shells in naval warfare in 1640'' is inconsistent with historical records documenting earlier use by other naval powers. Similarly, the assertion that ``London's population doubled between 1800 and 1820'' oversimplifies gradual demographic change; empirical population estimates for that period do not support a doubling.

In contrast, categories requiring precise technical knowledge, such as Mathematics (5.6\%) and Technology (9.4\%), show markedly lower inconsistency rates, suggesting that factual precision is better maintained in domains with more quantifiable information. Overall, these results indicate that Wikipedia's reliability varies across knowledge domains, with narrative-heavy subjects being particularly susceptible to inconsistencies.






\section{More Details on the Annotation Process for \dataset}\label{sec:annotation}

\subsection{Annotation Tool}
Figure~\ref{fig:dataset_construction} provides an overview of the dataset construction process.

\begin{figure*}[ht!]
\centering
\includegraphics[width=\linewidth]{assets/dataset_construction.pdf}
\caption{Overview of the \dataset~construction process: diverse passage sampling from Wikipedia's Vital Articles, adversarial claims collection using GPT-4o and a weak baseline filter, and human verification with detailed evidence analysis. Randomly sampled articles are used to estimate the prevalence of inconsistencies in the entire English Wikipedia.}
\label{fig:dataset_construction}
\end{figure*}

\paragraph{Filter to Balance Inconsistent Labels in the Dataset.}\label{sec:weak_baseline}
We implement a weak baseline to provide a permissive standard for inconsistency detection and filter out obviously consistent claims during \dataset construction. This baseline is a simplified version of the retrieve-and-verify system described in Section\ref{sec:experiment}, where the verifier outputs a binary inconsistency decision rather than a confidence score. We use GPT-4o mini~\citep{gpt4o} as the language model for these binary decisions.

\paragraph{Report Generation.}\label{sec:report_generation}
We develop a report generation system that produces a detailed report for each fact in the dataset. This system mirrors the setup in Section~\ref{sec:experiment} but replaces the verifier with a report generation stage. The report stage takes all retrieved evidence and the clarifications made by the tools agent and generates a two-sided analysis via two GPT-4o calls: one soliciting reasoning that the fact is inconsistent and another soliciting reasoning that it is consistent. This provides annotators with balanced information for final judgment. The final report shown to annotators includes both lines of reasoning and the agent's trace. See Figure~\ref{fig:annotation_tool_2} through~\ref{fig:annotation_tool_4} for illustrations.

\paragraph{Annotation Portal.} We built a web-based annotation platform to streamline the workflow. Annotators first check the extracted fact for any extraction issues, then review the detailed inconsistency analysis report, and finally assign a label of either ``consistent'' or ``inconsistent''. Screenshots of the interface are shown in Figure~\ref{fig:annotation_tool_1} through~\ref{fig:annotation_tool_4}.

\begin{figure*}[ht!]
\centering
\includegraphics[width=\textwidth]{assets/annotation_tool_1.png}
\caption{Screenshot 1 of the annotation tool showing the main interface for claim verification and inconsistency labeling.}
\label{fig:annotation_tool_1}
\end{figure*}

\begin{figure*}[ht!]
\centering
\includegraphics[width=\textwidth]{assets/annotation_tool_2.png}
\caption{Screenshot 2 of the annotation tool showing the main interface for claim verification and inconsistency labeling.}
\label{fig:annotation_tool_2}
\end{figure*}

\begin{figure*}[ht!]
\centering
\includegraphics[width=\textwidth]{assets/annotation_tool_3.png}
\caption{Screenshot 3 of the annotation tool showing the main interface for claim verification and inconsistency labeling.}
\label{fig:annotation_tool_3}
\end{figure*}

\begin{figure*}[ht!]
\centering
\includegraphics[width=\textwidth]{assets/annotation_tool_4.png}
\caption{Screenshot 4 of the annotation tool showing the main interface for claim verification and inconsistency labeling.}
\label{fig:annotation_tool_4}
\end{figure*}

\subsection{Annotators}
The dataset is annotated partly by the authors and partly by annotators recruited via Prolific~\cite{prolific}. To recruit external annotators, we conducted an initial qualification test using 10 randomly sampled facts from a subset previously labeled by the authors. Candidates were evaluated on both labeling accuracy and the quality of their written justifications. We selected the top 17 candidates who demonstrated strong analytical skills and high accuracy in identifying inconsistencies. Because inconsistency detection is nuanced, we required annotators to be native English speakers from the US or UK, hold a graduate degree (Masters or PhD), and maintain a Prolific approval rate of at least 95\%. The final pool of 17 annotators spent an average of 6.5 minutes evaluating each fact.

\subsection{Annotation Guidelines}
The task is inherently complex, which increases the risk of labeling errors. Determining whether a claim is inconsistent with a set of evidence is substantially more challenging than many standard annotation tasks. The definition of inconsistency can be context-dependent. For example, if a claim states a population of 4.8 million while the evidence reports 5 million, the case could be labeled ``inconsistent'' under exact matching or ``consistent'' if rounding is deemed acceptable. Likewise, comparing an imprecise expression such as ``a few years'' to a specific value is nontrivial, and the threshold for inconsistency is not always clear.

To mitigate ambiguity, we established explicit guidelines for such cases and instructed annotators to follow them closely.

\section{Dataset Details}

\subsection{Distribution of the Topics}\label{sec:dataset_analysis}
Figure~\ref{fig:topic_distribution} shows the distribution of topics covered by the facts.

\begin{figure*}[ht!]
\centering
\includegraphics[width=0.6\linewidth]{assets/category_distribution.pdf}
\caption{Distribution of topics across the \dataset dataset, showing the diversity of knowledge domains covered.}
\label{fig:topic_distribution}
\end{figure*}

\subsection{Dataset Examples}
\label{sec:dataset_examples}
Figure~\ref{tab:dataset_examples} presents detailed examples of claims with accompanying evidence from the dataset.

\begin{figure*}[htbp]
\centering
\includegraphics[width=\linewidth]{assets/example_table.pdf}
\caption{Examples of claims with evidence from the dataset and required competencies.}
\label{tab:dataset_examples}
\end{figure*}





\section{More Details on the User Study}
\label{sec:user_study}

\subsection{Browser Extension Implementation}

We implement the browser extension using JavaScript for the frontend and Python for the backend server. The frontend is a lightweight Chrome extension that injects content scripts to highlight potentially inconsistent claims on Wikipedia pages and provide a potential explanation for the inconsistency in the form of a side panel. For fact extraction, we use GPT-4o to parse Wikipedia page content into atomic claims, following prior work~\citep{semnani-etal-2023-wikichat,min-etal-2023-factscore}. The extension communicates with the backend via REST API endpoints. 

\begin{figure*}[ht!]
    \centering
    \includegraphics[width=\textwidth]{assets/chrome_extension_side_panel.png}
    \caption{The browser extension is implemented as a button in Wikipedia which the user can click to check for inconsistencies. When a claim is flagged as potentially inconsistent, we highlight the claim on the page and show a side panel with a detailed explanation of the inconsistency and links to the evidence documents.}
    \label{fig:browser_extension}
\end{figure*}

\subsection{Recruitment Details}

We recruited participants through by posting a research participation call on Meta-Wiki\footnote{\url{https://meta.wikimedia.org}}.
All recruited editors had made at least 200 edits and been active for over one year. The study protocol was approved by our institution's Institutional Review Board (IRB), and all participants provided informed consent before beginning the study.

\subsection{Perceived Usefulness and Qualitative Feecback}

In addition to Likert-scale ratings, we collect open-ended feedback about participants' experiences finding inconsistencies with and without our tool. Editors report that they employ diverse manual strategies, such as cross-checking linked articles, search and keyword search, reviewing talk pages, and validating references, but find these approaches time-consuming and cognitively cumbersome. They further express that they value the tool's accessibility for novice editors and utility for verifying AI-generated content. The main concerns center on processing speed, false positive rates, and interface design. These insights reveal a key tradeoff: while our system significantly reduces the editorial burden, its effectiveness ultimately depends on balancing detection sensitivity with efficiency and usability. See Table~\ref{tab:user_feedback} for examples of user feedback.

\begin{table*}[htbp]
\centering
\begin{tabular}{p{0.95\linewidth}}
\toprule
\textbf{Positive Feedback} \\
\midrule
The fact that I can simply enable the extension and, within a few minutes, see at a glance 126 inconsistencies is impressive. I also appreciate that it directly links me to where each inconsistency occurs in the article and provides a brief summary. Additionally, I like the feature that allows me to validate these inconsistencies by either accepting or rejecting them. \\
\midrule
It gave a lot of potential statements and claims that could be inconsistent with other information. One of the harder parts without using the tool was to find relevant articles that would include overlapping information. \\
\midrule
The tool is instructive, directive and straight to the point. It identifies inconsistency without rigorous means. \\
\midrule
Very good at finding discrepancies. \\
\midrule
That each result has an option to give feedback if the statement is actually inconsistent. \\
\midrule
Really helpful automation, especially for editors that may not be a subject-matter expert. Like that it could be used to fact-check AI-generated content if/when it comes to Wikipedia. \\
\midrule
It clearly demarked which statements were potentially inconsistent and with what, plus it allows for feedback on if the statements are in fact inconsistent. \\
\midrule
I really like the intent of the tool! I feel like it holds a lot of promise to help editors fact-check and correct inaccurate articles more quickly. I also love the UI - it's simple, clear, and easy to interpret. I especially appreciated the explanations offered in the panel, which even linked to helpful sources! That part was really impressive. I was able to detect one very clear inconsistency using this tool, and I was able to identify and validate it extremely quickly because of the helpful explanation and sources that the tool provided. \\
\midrule
\textbf{Areas for Improvement} \\
\midrule
I didn't like how often the tool was wrong. Even if, on paper, it would be better to highlight claims that the tool is unsure is consistent, in reality I personally felt very annoyed every time it was wrong. \\
\midrule
It assumed that Georgian, Renaissance Revival, etc. architecture were mutually exclusive with Palladian architecture, and that all the American buildings listed as having been influenced by Palladian architecture were inconsistent due to the architectural style listed in those articles. \\
\midrule
UI/UX improvements mostly. Would be nice if it was a native feature that could be turned on in Wikipedia and did not require downloading onto your computer. \\
\midrule
There wasn't much to dislike! I think that the model's precision could be improved, but the experience itself was really great. \\
\bottomrule
\end{tabular}
\caption{Qualitative User Feedback on \system}
\label{tab:user_feedback}
\end{table*}






\section{Implementation Details}
\label{sec:system_implementations}
We accessed OpenAI models via Azure OpenAI and accessed LLaMA models through Azure AI Services. For all experiments, greedy decoding (i.e. temperature 0) is used.
All numbers are the result of a single run.

\subsection{Implementation of Explain and Clarify Tools}
\label{sec:tools_implementation}

The implementation of clarification tools in our system enables the verification agent to gather additional information when evaluating a claim. 

\begin{itemize}
    \item \action{explain} prompts the LLM to provide background information about the topic query. As shown in Figure~\ref{lst:generate_background_info}, this action instructs the LLM to synthesize its existing knowledge about the topic in relation to the claim being evaluated.
    \item \action{clarify} disambiguates entities by first retrieving relevant information and then using the language model to explain differences based on that retrieved information. The number of retrieved documents per each clarify is 10. As shown in Figure~\ref{lst:generate_entity_report} and~\ref{lst:generate_entity_report_2}, we prompt the language model to analyze the retrieved results to resolve ambiguities by explain the differences.
\end{itemize}

\subsection{Prompts for Systems}

We list all the prompts used for implementing the agent system below. Figure ~\ref{lst:claim_extraction} is our prompt for extracting atomic facts from Wikipedia articles, whereas Figures ~\ref{lst:generate_background_info}, ~\ref{lst:generate_entity_report}, ~\ref{lst:generate_entity_report_2}, ~\ref{lst:generate_entity_report_3},~\ref{lst:verifier}, ~\ref{lst:verifier_2}, ~\ref{lst:verifier_3}, and ~\ref{lst:controller} are the tools available to the \system agent.
In each prompt, \texttt{\# input} and \texttt{\# output} denote the boundaries of few-shot examples used, if any.

\begin{figure*}[htb]
\begin{lstlisting}
# instruction
You are an expert fact extractor tasked with identifying and listing atomic facts from a given text. Your goal is to produce a comprehensive list of facts that are explicitly stated or directly inferrable from the provided information.

Instructions:
1. Read the title and text carefully.
2. Extract all atomic facts from the information provided. An atomic fact is a single, indivisible piece of information that cannot be broken down further without losing its meaning or accuracy.
3. Include only facts that are explicitly stated or can be directly and unambiguously inferred from the text.
4. Do not add any external knowledge or assumptions not present in the given information.
5. Ensure that each fact is self-contained and can be independently fact-checked.

Before providing your final list of facts, break down your fact extraction process in <fact_extraction_process> tags. This will help ensure a thorough and accurate extraction of facts.

In your fact extraction process, follow these steps:
1. Identify key topics or themes from the title and text.
2. For each topic/theme, list explicit facts from the text.
3. Consider potential inferences that can be directly drawn from the explicit facts, and evaluate their validity.
4. Evaluate each fact (explicit and inferred) for atomicity and self-containment.
5. Categorize facts by topic/theme.
6. Cross-reference each fact with the original text to ensure accuracy.
7. Review the list to ensure no redundant or overlapping facts are included.

After your analysis, provide your final list of facts, with each fact on a new line.

Example output structure:

<fact_extraction_process>
[Your detailed fact extraction process, following the steps outlined above]
</fact_extraction_process>

<facts>
[Fact 1]
[Fact 2]
[Fact 3]
...
</facts>

# input
Here is the title and text you need to analyze:

<title>
{{ full_title }}
</title>

<text>
{{ text }}
</text>
\end{lstlisting}

\caption{Fact Extraction Prompt}
\label{lst:claim_extraction}
\end{figure*}


\begin{figure*}[htb]
\begin{lstlisting}
# instruction
You will be given a topic, and a Wikipedia passage where the topic is mentioned. Your task is to write a self-contained paragraph explaining technical or domain-specific terms in the topic. Your goal is to provide background information on the given topic for people who are unfamiliar with it. If a term, event or concept in the topic has multiple interpretations or meanings, list all plausible ones.

# input
Topic: Infanta Amalia
Wikipedia article: Infanta Amalia of Spain
Infanta Amalia of Spain (Spanish: Amalia de Borbon y Borbon-Dos Sicilias; 12 October 1834 - 27 August 1905) was the youngest daughter of Infante Francisco de Paula of Spain. Her eldest brother, Francisco de Asis, married Queen Isabella II of Spain, who was Amalia's first cousin.

# output
"Infanta Amalia" refers to a title and name in Spanish and Portuguese contexts. "Infanta" is a title used in Spain and Portugal for the daughters of a monarch who are not heir apparent, similar to "princess" in English. "Amalia" is a given name. Therefore, "Infanta Amalia" would refer to a princess named Amalia within a royal family in Spain or Portugal.

# input
Topic: The Great Gatsby
Wikipedia article: The Great Gatsby
It was also performed in the summer of 2012 at the Aspen Music Festival and School. It was performed at Seagle Festival in Schroon Lake, NY in the summer of 2018.

# output
"The Great Gatsby" here likely to a musical adaptation, play, opera, or other performance based on the novel "The Great Gatsby" by F. Scott Fitzgerald. The novel is a classic work of American literature published in 1925. The performances mentioned in the passage are likely adaptations of the novel for the stage or other artistic mediums.

# input
Topic: {{ topic }}
Wikipedia article: {{ claim.context_block.full_title }}
{{ claim.context_block.content }}
\end{lstlisting}

\caption{Generate Background Information Prompt}
\label{lst:generate_background_info}
\end{figure*}




\begin{figure*}[htb]
\begin{lstlisting}
# instruction
You will be given an entity and a Wikipedia paragraph where it is mentioned.
You will also be provided with a list of search results that may contain information about the entity, and other similar entities.
Your task is to write a self-contained paragraph explaining the differences between entities with similar names in the search results.
Entities with similar names might lead to confusion, and the goal here is to disambiguate them. Pay attention to the following:

- People with the same last name, but different first names. Or People with the same name but different professions or time periods.
- Events with the same name but different years or locations. For example, "The Olympics" could refer to the winter or summer games, or games held in different years.
- Organizations with similar names but different purposes or locations.
- etc.

# input
Entity: members of the royal family of Spain named Amalia

[1] Title: Infanta Maria Amalia of Spain
Maria Amalia, Infanta of Spain (9 January 1779 in Madrid - 22 July 1798 in Madrid), was a Spanish princess. She was a daughter of King Charles IV of Spain, in 1795, she married her uncle Infante Antonio Pascual of Spain.

[2] Title: Infanta Amalia of Spain > Childhood
She was born at the royal Palace of Madrid on 12 October 1834 as the eleventh child and sixth daughter of Infante Francisco de Paula of Spain, younger brother of King Fernando VII of Spain, and his wife, Princess Luisa Carlota of Bourbon-Two Sicilies. Infanta Amalia's mother was the niece of her father since her maternal grandmother, Infanta Maria Isabella of Spain, was the elder sister of Infante Francisco de Paula.

[3] Title: Infanta Maria Amalia of Spain > Early life
Born at the Royal Palace of El Pardo, Maria Amalia was the second surviving daughter of King Carlos IV of Spain (1748-1819) and his wife Maria Luisa of Parma (1751-1819), a granddaughter of Louis XV of France.

[4] Title: Infanta Amalia of Spain
Infanta Amalia of Spain (Spanish: Amalia de Borbon y Borbon-Dos Sicilias; 12 October 1834 - 27 August 1905) was the youngest daughter of Infante Francisco de Paula of Spain. Her eldest brother, Francisco de Asis married Queen Isabella II of Spain, who was Amalia's first cousin. She was one of only two of five sisters who made a royal marriage. In 1865 she married Prince Adalbert of Bavaria, a son of King Ludwig I of Bavaria. Upon her marriage she moved to Munich, where she spent the rest of her life. However she remained attached to her native country and was instrumental in arranging the marriage of her eldest son Prince Ludwig Ferdinand of Bavaria with her niece Infanta Paz of Spain.

[5] Title: Infanta Amalia of Spain > Later life and death
Although Infanta Amalia lived for the rest of her life in Munich, she remained attached to her native country. She visited Spain often and her eldest son Prince Ludwig Ferdinand of Bavaria was born at the royal palace of Madrid. She spent the winters at the residence of Munich and the summers at Nymphenburg Palace. Her husband died in 1875; Amalia outlived him by thirty years. Amalia maintained her affiliation with Spain in the next generation. All of her five children spoke Spanish fluently and she encouraged her son Ludwig Ferdinand to marry her niece and goddaughter Infanta Maria de la Paz of Spain. The couple married in 1883.

\end{lstlisting}

\caption{Generate Entity Report Prompt}
\label{lst:generate_entity_report}
\end{figure*}






\begin{figure*}[htb]
\begin{lstlisting}
# output
There are two entities with similar names.
    1. "Infanta Amalia of Spain": Infanta Amalia of Spain (Spanish: Amalia de Borbon y Borbon-Dos Sicilias; 12 October 1834 - 27 August 1905) was the youngest daughter of Infante Francisco de Paula of Spain.
    2. "Infanta Maria Amalia of Spain": Maria Amalia, Infanta of Spain (9 January 1779 in Madrid - 22 July 1798 in Madrid), was a Spanish princess. She was a daughter of King Charles IV of Spain, in 1795, she married her uncle Infante Antonio Pascual of Spain.


These two individuals seem to be separate entities, but may be relatives.

# input
Entity: Antoine Emile Henry Labeyrie

[1] Title: Antoine Emile Henry Labeyrie
Antoine Emile Henry Labeyrie (born 12 May 1943) is a French astronomer, who held the Observational astrophysics chair at the College de France between 1991 and 2014, where he is currently professor emeritus. He is working with the Hypertelescope Lise association, which aims to develop an extremely large astronomical interferometer with spherical geometry that might theoretically show features on Earth-like worlds around other suns, as its president. He is a member of the French Academy of Sciences in the Sciences of the Universe (sciences de l'univers) section. Between 1995 and 1999 he was director of the Haute-Provence Observatory.

[2] Title: Galluis > Notable residents
Antoine-Germain Labarraque (1777 - 1850) was a French chemist and pharmacist, notable for formulating and finding important uses for "Eau de Labarraque" or "Labarraque\'s solution", a solution of sodium hypochlorite widely used as a disinfectant and deodoriser. He died in Gallius on 9 December 1850.

[3] Title: Antoine Lavoisier
Antoine-Laurent de Lavoisier (26 August 1743 - 8 May 1794), also Antoine Lavoisier after the French Revolution, was a French nobleman and chemist who was central to the 18th-century chemical revolution and who had a large influence on both the history of chemistry and the history of biology.

[4] Title: Antoine Germain Labarraque
Antoine Germain Labarraque (28 March 1777 - 9 December 1850) was a French chemist and pharmacist, notable for formulating and finding important uses for "Eau de Labarraque" or "Labarraque\'s solution", a solution of sodium hypochlorite widely used as a disinfectant and deodoriser.

[5] Title: Antoine Germain Labarraque
| Antoine Germain Labarraque | |
| --- | --- |
| Portrait of Labarraque | |
| Born | (1777-03-28) Oloron-Sainte-Marie, Pyrenees-Atlantiques, France |
| Died | 9 December 1850 (1850-12-09) (aged 73)near Paris, France |
| Nationality | French |
| Education | College of Pharmacy, Paris |
| Occupation(s) | chemist and pharmacist |
| Known for | using sodium hypochlorite as a disinfectant and deodoriser |
| Parents | * Francois Labarraque (father) * Christine Sousbielle (mother) |
\end{lstlisting}

\caption{Generate Entity Report Prompt (Continued)} \label{lst:generate_entity_report_2}
\end{figure*}

\begin{figure*}[htb]
\begin{lstlisting}
# output
There are multiple notable French scientists with similar names beginning with "Antoine".
  1. Antoine Emile Henry Labeyrie (born 1943) is a French astronomer and professor emeritus who held the Observational astrophysics chair at the College de France.
  2. Antoine-Laurent de Lavoisier (1743-1794) was a French nobleman and chemist central to the 18th
  3. Antoine Germain Labarraque (1777-1850) was a French chemist and pharmacist known for developing "Labarraque's solution," a sodium hypochlorite disinfectant.
While these individuals share similar first names and French nationality, they worked in different fields and time periods.

# input
Entity: {{ entity_name }}
Original article: {{ claim.context_block.full_title }}
{{ claim.context_block.content }}

{{ search_results.to_string() }}
\end{lstlisting}

\caption{Generate Entity Report Prompt (Continued)} \label{lst:generate_entity_report_3}
\end{figure*}


\begin{figure*}[htb]
\begin{lstlisting}
# instruction
Determine if a claim extracted from a Wikipedia paragraph is inconsistent with any of the provided documents. A claim is deemed inconsistent when at least one document contains information that directly contradicts it. If no such contradiction exists - even when the documents do not explicitly support the claim - the claim is considered consistent.

Step-by-Step Instructions:

1. Identify the Claim
Definition: A brief statement directly extracted from a Wikipedia paragraph.
Note: The full meaning of the claim might require context provided by the original paragraph.

2. Review the Documents
Definition: Passages, tables, or pieces of text retrieved from Wikipedia.
Task: Ignore documents that are clearly irrelevant to the claim.
Focus on finding any document that might contain information in clear conflict with the claim.

3. Consider Clarifications

Definition: Additional background information provided to clarify ambiguous terms or entities.
Task: Use clarifications to distinguish between similar or similarly named entities.
Important: Do not use clarifications to support or contradict the claim directly - they serve only to clear up ambiguities.

4. Assess for Inconsistencies

Definition of Inconsistency:
The claim is inconsistent if at least one document provides information that contradicts it.
Conversely, if no document provides conflicting information, the claim is considered consistent.
Measurement: Assign an inconsistency score between 0 (fully consistent) and 1 (completely inconsistent).
Intermediate scores indicate varying degrees of uncertainty or partial conflict.

5. Common Scenarios & Examples

Example 1: Clear Inconsistency

Claim: "The capital of Thailand is Bangkok."
Document: "The capital of Thailand is Phuket."
Reasoning: A country typically has one capital. The document contradicts the claim by listing a different city, yielding a high inconsistency score (e.g., 0.8-0.9).

Example 2: Apparent Inconsistency Resolved by Entity Equivalence (Minor Inconsistency)

Claim: "The capital of Thailand is Bangkok."
Document: "The capital of Thailand is Krung Thep Maha Nakhon."
Additional Background: It is widely accepted that Bangkok and Krung Thep Maha Nakhon refer to the same city.
Reasoning: Although the names differ, they reference the same location; thus, the claim is largely consistent (e.g., inconsistency score around 0.2-0.4).
Note: If an explicit clarification were provided stating the equivalence, the score would be 0.

Example 3: Misplaced Terms Causing Inconsistency

Claim: "The capital of Thailand is Bangkok."
Document: "The capital of Bangkok is Thailand."
\end{lstlisting}

\caption{Verifier Prompt} \label{lst:verifier}
\end{figure*}


\begin{figure*}[htb]
\begin{lstlisting}
Reasoning: The document seems to mix up entities by stating that Bangkok is a country. With no supporting evidence that this is a mere typo or misinterpretation, the conflict earns a high inconsistency score (e.g., around 0.9).

Example 4: Inconsistent Translational Variants

Claim: "The 'Song is Universal' won the Best Modern Rock Song award at the 2010 Korean Music Awards."
Document: "The Best Modern Rock Song award at the 2010 Korean Music Awards was given for 'Universal Song.'"
Additional Clarification: "Bangkok only refers to a city in Thailand, not elsewhere." (Not directly applicable here but shows how clarifications work.)
Reasoning: Although the song likely is the same, the differing English translations ("Song is Universal" vs. "Universal Song") introduce an inconsistency, resulting in a moderately high inconsistency score (e.g., around 0.8).

Example 5: No Conflict (Consistency)

Claim: "Stress is harmful to health, as mentioned in the medical literature."
Document: "Stress is necessary for growth and development, pushing limits, enhancing learning, and building resilience."
Reasoning: The document discusses the beneficial aspects of acute or eustress compared to chronic stress, which is what the claim addresses. Since these are two different perspectives on stress, there is no contradiction - the claim is consistent (inconsistency score 0).

Final Decision

After review, provide the inconsistency score for the claim:
0: Fully consistent; no document contradicts the claim.
Between 0 and 1: Partial or potential inconsistencies.
1: Fully inconsistent; at least one document directly contradicts the claim.

# input
<claim>
Title: {{ claim.context_block.full_title }}
{{ claim.context_block.content }}

You should only focus on the aspect of this paragraph related to: "{{ claim.claim_text }}"
</claim>

Read the following clarifications about the claim:
<clarifications>
{%
[{{ loop.index }}] {{ clarification }}

{%
</clarifications>

Now, read through the documents below and look for any information that conflicts with the claim:
<documents>
{%
[{{ loop.index }}] {{ document }}

{%
</documents>
\end{lstlisting}

\caption{Verifier Prompt (Continued)} \label{lst:verifier_2}
\end{figure*}


\begin{figure*}[htb]
\begin{lstlisting}
Now, you need to analyze the documents and clarifications to determine an inconsistency score that represents your confidence that the claim is inconsistent with the documents.

First, rephrase the claim to be more specific by:
1. Incorporating context from the Wikipedia article title and content
2. Preserving the original meaning, but making corrections if the claim appears to be misrepresented or incorrectly paraphrased from the claim's context.

<claim_with_context>
[Provide the claim that incorporates the context from the Wikipedia article title and content here.]
</claim_with_context>

Based on the claim with context, present your full analysis and arguments:

<analysis>
[Provide a detailed analysis by:
1. Carefully examining the claim and documents for any contradictions or inconsistencies, look through examples above if there is concept similar to your case
2. Highlighting specific documents where information directly conflicts with the claim
3. Making sure that these documents are relevant to the claim. Some documents may contain the same entities as the claim, but they are not relevant to the claim, given the context
4. Exploring multiple interpretations of the claim's meaning and implications
5. Considering edge cases and ambiguities that could affect the analysis
6. Referencing relevant examples from above (translations, time-related issues, ordering) to strengthen your reasoning
7. Explaining your confidence level in identifying any inconsistencies found]

</analysis>

Based on your analysis, provide an inconsistency score from 0 to 1, where:
- 0 indicates the claim is completely consistent with all of the documents
- 1 indicates the claim is completely inconsistent with at least one of the documents
- Values between 0 and 1 represent varying degrees of uncertainty

<inconsistency_score>
[A single float from 0 to 1]
</inconsistency_score>
\end{lstlisting}

\caption{Verifier Prompt (Continued)} \label{lst:verifier_3}
\end{figure*}


\begin{figure*}[htb]
\begin{lstlisting}
# instruction
You will be given a "claim" statement extracted from a Wikipedia paragraph.
Your task is to conduct a thorough investigation on the entire Wikipedia (except the article where the claim comes from) to find any factual inconsistencies with this claim.
As you conduct your investigation, you may come across articles that support the claim. However, you should continue searching for inconsistencies that might exist in other places. Inconsistencies might appear in subtle or indirect ways.

You will conduct your investigation in multiple steps. At each step, you should think about the information you have gathered so far, and choose one of the following actions based on it:

- `explain(topic: str) -> str`: Use this action to understand the basics of a specific term or concept you encounter, for example a technical term or the rules of a sport.

- `clarify_entity(entity_name_and_description: str) -> str`: Use this action to get a report on an entity (person, organization, event etc.) to clarify other entities with similar names. This will help you properly differentiate similar-sounding entities when researching inconsistencies. For example, clarify_entity("WW III wrestling event") will explain all potential events with similar names, or the same event in different years.

- `search_wikipedia_outside_claim_article(question: str) -> list`: Use this action to explore Wikipedia.

- `report_inconsistency(evidence: str)`: If at any point you are certain that you have found an inconsistency, use this action to report it. Evidence should be a short sentence that describes the inconsistency. Once you report an inconsistency, a human will review it and provide feedback.

# input
Here is the claim to find inconsistencies with:
{{ claim.claim_text }}

Here is more context about the claim for your reference:
Title: {{ claim.context_block.full_title }}
{{ claim.context_block.content }}

{%
Actions you have taken so far:

{{ action_history}}
{%
\end{lstlisting}

\caption{Controller Prompt}
\label{lst:controller}
\end{figure*}








































